
\section{Background}

\subsection{Truncated SVD}

For any matrix $W \in \mathbb{R}^{m \times n}$, the SVD is:

\[
W = U \Sigma V^\top
\tag{3.1}
\]

where $U \in \mathbb{R}^{m \times m}$, $V \in \mathbb{R}^{n \times n}$ are 
orthogonal, and $\Sigma$ is diagonal with non-negative entries 
$\sigma_1 \geq \sigma_2 \geq \dots \geq 0$.

\textbf{Truncated SVD} keeps the top-$r$ singular values:

\[
W_r = U_{:, :r} \Sigma_{:r, :r} V_{:r, :}^\top
\tag{3.2}
\]

with error $\|W - W_r\|_F = \sqrt{\sum_{i > r} \sigma_i^2}$.

\textbf{Symmetric factorization.} For numerical stability, we distribute 
$\Sigma^{1/2}$ between the two factors:

\[
U_r = U_{:, :r} \Sigma_{:r, :r}^{1/2}, \quad
V_r = \Sigma_{:r, :r}^{1/2} V_{:r, :}^\top
\tag{3.3}
\]

so that $W_r = U_r V_r$, and the two factors have balanced magnitude.

\subsection{Input Activation Covariance}

For a linear layer $y = Wx + b$ with input $x \sim \mathcal{D}$, define:

\[
C_X = \mathbb{E}_{x \sim \mathcal{D}}\!\left[x x^\top\right] \in \mathbb{R}^{d_\text{in} \times d_\text{in}}
\tag{3.4}
\]

This is a symmetric positive semi-definite matrix. Its eigendecomposition:

\[
C_X = Q \Lambda Q^\top, \qquad
\Lambda = \mathrm{diag}(\lambda_1, \dots, \lambda_{d_\text{in}})
\tag{3.5}
\]

gives orthogonal principal directions (columns of $Q$) and corresponding 
variances ($\lambda_i$).

\textbf{Diagonal case.} If input dimensions are independent, $C_X$ is diagonal and 
$C_X = \mathrm{diag}(F_1, \dots, F_{d_\text{in}})$ where $F_i = \mathbb{E}[x_i^2]$. 
This is the "Fisher diagonal" used in A3.

\textbf{General case.} Real activations are \textbf{not} independent -- attention 
residual streams and FFN hidden states exhibit strong correlations. AASVD 
captures this via the full matrix $C_X$.

\subsection{Matrix Square Root}

For a positive semi-definite matrix $C_X = Q \Lambda Q^\top$, the matrix 
square root is:

\[
C_X^{1/2} = Q \Lambda^{1/2} Q^\top, \qquad
\Lambda^{1/2} = \mathrm{diag}(\sqrt{\lambda_1}, \dots, \sqrt{\lambda_{d_\text{in}}})
\tag{3.6}
\]

The inverse square root is defined similarly with $1/\sqrt{\lambda_i}$.

\textbf{Regularization.} For numerical stability, we regularize:

\[
\lambda_i \leftarrow \max(\lambda_i, \epsilon_\text{reg}), \qquad
\epsilon_\text{reg} = 10^{-6} \max_j \lambda_j + 10^{-8}
\tag{3.7}
\]

\subsection{Exponential Moving Average}

Given a stream of covariance estimates $C_1, C_2, \dots, C_K$ (one per 
calibration batch), the EMA is:

\[
\hat{C}^{(k)} = \beta \hat{C}^{(k-1)} + (1 - \beta) C_k, \qquad
\hat{C}^{(0)} = C_1
\tag{3.8}
\]

with decay $\beta \in [0, 1]$. For $\beta = 0.9$, recent batches get weight 
$(1 - \beta) = 0.1$, giving an effective window of ~10 batches.

\textbf{Why EMA?} Straight mean weighting is $\beta \to 0$. EMA gives 
\textbf{recency bias}, useful when the calibration data comes from a distribution 
with drift. For stationary distributions, both give the same expected value.

\textbf{Memory.} EMA stores only $\hat{C}^{(k)}$, a $d \times d$ matrix. Total 
memory is $\mathcal{O}(N_l d^2)$, independent of the number of calibration 
batches.

\subsection{LoRA Recovery}

LoRA (Hu et al., 2021) parameterizes weight updates as:

\[
\Delta W = \frac{\alpha}{r_L} B A
\tag{3.9}
\]

with $A \in \mathbb{R}^{r_L \times d_\text{in}}$, $B \in \mathbb{R}^{d_\text{out} \times r_L}$, 
$B = 0$ at initialization. Only $A, B$ are trained.

\textbf{In AASVD.} We add LoRA adapters to the factorized layer:

\[
y = (x V^\top) U^\top + (x A^\top) B^\top \cdot s + b
\tag{3.10}
\]

where $s = \alpha / r_L$. The frozen $U, V$ provide the bulk; LoRA provides 
recovery capacity.

\subsection{Frobenius Reconstruction Error}

For validation, we use the relative Frobenius error:

\[
\text{err}(W, U, V) = \frac{\|W - UV\|_F}{\|W\|_F + \epsilon}
\tag{3.11}
\]

This measures how much of the weight matrix's energy is lost. Thresholds 
are typically $0.10$ (tight) to $0.50$ (permissive).

\subsection{Notation Summary}
